\chapter{Richiami}
La supergravità e la corrispondenza AdS/CFT costituiscono due importanti sviluppi nello studio della gravità quantistica, della supersimmetria e della teoria delle stringhe.

La relatività generale descrive la gravità come una teoria geometrica dello spaziotempo, mentre la teoria quantistica dei campi fornisce il quadro teorico per la descrizione delle interazioni fondamentali a livello quantistico. La costruzione di una teoria quantistica della gravità presenta tuttavia difficoltà profonde, che motivano la ricerca di strutture teoriche più generali rispetto alla teoria quantistica dei campi ordinaria.

La teoria delle stringhe fornisce una possibile estensione di questo quadro, nella quale le particelle fondamentali sono sostituite da oggetti estesi e la gravità emerge naturalmente dallo spettro della teoria. Nel limite di basse energie, le teorie di stringa sono descritte da teorie efficaci di supergravità, nelle quali la supersimmetria è realizzata localmente e la dinamica gravitazionale è accompagnata da ulteriori gradi di libertà. La supersimmetria introduce una relazione tra bosoni e fermioni e impone forti vincoli sulla struttura delle teorie, sulle loro simmetrie e sulle loro possibili interazioni.

La corrispondenza AdS/CFT propone inoltre una relazione profonda tra una teoria gravitazionale e una teoria quantistica dei campi, mostrando come descrizioni apparentemente diverse possano risultare equivalenti. Questa dualità fornisce un nuovo punto di vista sulla gravità quantistica e sulle teorie fortemente accoppiate, e costituisce uno dei principali strumenti teorici per esplorare la struttura non perturbativa delle teorie fondamentali.

\begin{notz}(unità naturali)
Nel seguito si opererà sempre in unità naturali, ponendo cioè 
\begin{equation}
    c = \hbar = k_B = 1;
\end{equation}
tuttavia, non in tutte le circostanze si pone $G = 1$, la scelta dipenderà dal contesto.
\end{notz}
\begin{notz}(notazione di Einstein)
    Nel seguito si adotterà, a meno che non sia specificato diversamente, la notazione di Einstein per la sommatoria.
\end{notz}

\section{Relatività speciale e teoria dei campi}
Anche se nel seguito si lavorerà su varietà di natura generale, come primo passo si richiama la varietà associata alla spaziotempo in relatività speciale. Considerando il caso $D$-dimensionale, tale varietà è una varietà Riemanniana in quanto è possibile definire una metrica, in particolare quella data da
\begin{equation}
\label{eq: metrica di Minkowski in D dimensioni}
    \eta_{\mu\nu} = \diag{(-1,\underbrace{\dots}_{D - 1\,\mathrm{volte}},1)},
\end{equation}
la cui inversa è definita come $\eta_{\mu\nu}^{-1} \equiv \eta^{\mu\nu}$, la quale coincide con la metrica $\eta_{\mu\nu}$. Tuttavia, per una varietà generica la metrica inversa non coincide necessariamente con la sua controparte. Ad ogni modo, sulla varietà dello spaziotempo si definiscono i vettori $D$-dimensionali come
\begin{equation}
    x^\mu = (t,x^i),\quad \mu = 0, \dots, D - 1,\,i = 1, \dots, D - 1.
\end{equation}
\begin{notz}(indice di Lorentz)
    Nel seguito, se non si incorre in ambiguità, si potrà omettere l'indice di Lorentz $\mu$ dei vettori $x^\mu$, permettendo così di alleggerire la notazione ponendo $x^\mu \equiv x$.
\end{notz}
Lo spazio che si sta descrivendo in questo modo corrisponde al prodotto $\bb{R} \ot \cdots \ot \bb{R}$ dato che le coordinate dei vettori $x^\mu$ possono variare sull'intervallo $(-\infty,\infty)$. Tramite la metrica \eqref{eq: metrica di Minkowski in D dimensioni} si può definire la norma di un vettore come
\begin{equation}
\label{eq: norma su spazio di Minkowski in D dimensioni}
    ||x^\mu||^2 \equiv x^\mu\eta_{\mu\nu}x^\nu = -t^2 + (x^i)^2 = x^\mu x_\mu.
\end{equation}
Si noti che in questo modo è possibile definire il vettore $x_\mu$ con l'indice in basso, in particolare
\begin{equation}
    x_\mu \equiv \eta_{\mu\nu}x^\nu;
\end{equation}
si noti quindi che questa non è una definizione posta a priori. Una nozione particolarmente importante in Relatività speciale è l'elemento di linea $ds^2$ definito come
\begin{equation}
    ds^2 = \eta_{\mu\nu}dx^\mu dx^\nu = -dt^2 + dx_idx^i,
\end{equation}
dove formalmente $dx^\mu dx^\nu$ corrisponde al prodotto $dx^\mu \ot dx^\nu$. Lo spazio così definito e avente queste proprietà è detto ``spazio di Minkowski'' in $D$ dimensioni, denotato con $\bb{R}^{1,D - 1}$.Su questo spazio agiscono delle trasformazioni lineari della forma $x^\mu \to x'^\mu \equiv \Lambda^\mu_\nu x^\nu$. Tra tutte queste trasformazioni sono interessanti quelle che lasciano la norma \eqref{eq: norma su spazio di Minkowski in D dimensioni} invariata, cioè tali che $||x^\mu||^2 = ||x'^\mu||^2$. Queste particolari trasformazioni formano un gruppo di Lie indicato con $\rem{O}(1,D - 1)$, dove, in forma matriciale, sono tali da soddisfare la relazione
\begin{equation}
    \Lambda^T\eta\Lambda = \eta,\quad \Lambda \in \rem{O}(1,D - 1).
\end{equation}
Se si considerano poi tra le trasformazioni $\Lambda$ appartenenti al gruppo $O(1,D - 1)$ quelle che hanno $\det{\Lambda} = 1$, si può dimostrare che queste formano un sottogruppo di Lie indicato con $SO(1,D - 1)$ e detto ``gruppo di Lorentz''. 

Gli oggetti con cui si avrà principalmente a che fare saranno i campi, tipicamente indicati con $\phi(x)$. In particolare, si farà uso della formulazione lagrangiana per una descrizione esplicita delle simmetrie. Il primo esempio di campo è il campo scalare, il quale è definito tramite la trasformazione 
\begin{equation}
    \phi(x) \to \phi'(x') \equiv \phi(x') = \phi(\Lambda x).
\end{equation}
Un altro esempio è il campo vettoriale, definito invece tramite la trasformazione
\begin{equation}
    V^\mu(x) \to V'^\mu(x') \equiv (\Lambda^{-1})^\mu_\nu V^\nu(\Lambda x).
\end{equation}
Un terzo esempio, molto importante, è il campo spinoriale, di cui però si tratterà meglio in seguito. 

In generale, l'elemento fondamentale per formulare una teoria è l'azione $S$, L'azione tipicamente può dipendere dai campi $\phi$ e dalle sue derivate $\p_\mu \phi$, contenute nella densità di lagrangiana $\Lag$. La definizione in quattro dimensioni, semplicemente estendibile in $D$ dimensioni, è 
\begin{equation}
\label{eq: definizione dell'azione in 4 dimensioni}
    S = \int d^4x\,\Lag(\phi,\p_\mu \phi).
\end{equation}
Se si compie una variazione arbitraria $\phi \to \phi + \delta\phi$, l'azione varia secondo
\begin{equation}
    \delta S = \int d^4x\bigg[\frac{\p\Lag}{\p\phi}\delta\phi - \p_\mu\bigg(\frac{\p\Lag}{\p(\p_\mu\phi)}\bigg)\delta\phi + \p_\mu\bigg(\frac{\p\Lag}{\p(\p_\mu\phi)}\delta\phi\bigg)\bigg].
\end{equation}
Imponendo adesso che la variazione dei campi calcolata sul bordo $\p$ della varietà si nulla, ossia che $\delta\phi|_\p = 0$, la variazione dell'azione si riduce a
\begin{equation}
    \delta S = \int d^4x\bigg[\frac{\p\Lag}{\p\phi}\delta\phi - \p_\mu\bigg(\frac{\p\Lag}{\p(\p_\mu\phi)}\bigg)\delta\phi\bigg].
\end{equation}
\begin{obs}
    Ci si potrebbe domandare cosa si intenda per bordo della varietà se tale varietà è infinita. Tipicamente, si intende il limite per $x \to \infty$ della quantità in questione
\end{obs}
Imponendo il principio di minima azione, si trova che $\delta S = 0$ se e solo se $\Lag$ soddisfa le equazioni di Eulero-Lagrange
\begin{equation}
\label{eq: equazioni di Eulero-Lagrange}
    \frac{\p\Lag}{\p\phi} - \p_\mu\bigg(\frac{\p\Lag}{\p(p_\mu\phi)}\bigg) = 0,
\end{equation}
ovvero le equazioni del moto.
\begin{obs}
    Le equazioni \eqref{eq: equazioni di Eulero-Lagrange} sono corrispondenti a quella che si deriva in Meccanica classica
    \begin{equation}
        \frac{\p L}{\p q} - \frac{d}{dt}\bigg(\frac{\p L}{\p \dot{q}}\bigg) = 0,
    \end{equation}
    dove $\Lag$ corrisponde ad $L$ e $\phi$ corrisponde a $q$.
\end{obs}

Per quanto riguarda le simmetrie, queste sono codificate dal teorema di Noether. L'idea è che ad ogni simmetria, ossia ad ogni legge di conservazione, corrisponde una quantità conservata. In generale, una simmetria è una trasformazione che lascia invariata l'azione, ossia tale che 
\begin{equation}
    S[\phi] = S[\phi'].
\end{equation}
Per definizione dell'azione \eqref{eq: definizione dell'azione in 4 dimensioni}, questa è invariante sotto trasformazioni della lagrangiana tramite una derivata totale, cioè per $\Lag \to \Lag' = \Lag + \epsilon\p_\mu k^\mu$. Per una trasformazione generale $\Lag \to \Lag' = \Lag + \epsilon\Delta\Lag$, si ha che 
\begin{equation}
    \begin{split}
        \epsilon\Delta\Lag &= \epsilon\p_\mu\bigg(\frac{\p\Lag}{\p(\p_\mu\phi)}\Delta\phi\bigg) - \epsilon\p_\mu\bigg(\frac{\p\Lag}{\p(\p_\mu\phi)}\bigg)\Delta\phi + \epsilon\frac{\p\Lag}{\p\phi}\Delta\phi \\
        &= \epsilon\p_\mu\bigg(\frac{\p\Lag}{\p(\p_\mu\phi)}\Delta\phi\bigg),
    \end{split}
\end{equation}
dove il campo trasforma come $\phi \to \phi' = \phi + \epsilon\Delta\phi$, la quale non è una trasformazione arbitraria, ma deve essere esplicitata per ogni caso. Ponendo ora che $\Delta\Lag = \p_\mu k^\mu$ e usando le equazione di Eulero-Lagrange \eqref{eq: equazioni di Eulero-Lagrange}, si definisce la corrente di Noether
\begin{equation}
    J^\mu := -k^\mu + \frac{\p\Lag}{\p(\p_\mu\phi)}\Delta\phi,
\end{equation}
la quale soddisfa la legge di conservazione 
\begin{equation}
\label{eq: legge di conservazione di Noether}
    \p_\mu J^\mu = 0.
\end{equation}
Dalla corrente di Noether si definisce la carica di Noether
\begin{equation}
    Q := \int d^3x\,J^0,
\end{equation}
da cui segue che 
\begin{equation}
    \frac{dQ}{dt} = 0.
\end{equation}

Si considera ora una traslazione nello spaziotempo 4-dimensionale $x^\mu \to x'^\mu = x^\mu + a^\mu$, che è una trasformazione appartenente al gruppo di Poincaré. Sotto questa trasformazione, un campo scalare trasforma come $\phi(x) \to \phi(x + a)$, dove $\phi(x + a)$ può essere espresso tramite uno sviluppo in serie di Taylor
\begin{equation}
    \phi(x + a) = \phi(x) + a^\mu\p_\mu\phi(x) + \cdots,
\end{equation}
dove si può identificare $\epsilon \equiv a^\mu$ e $\Delta\phi \equiv \p_\mu\phi(x)$. Pertanto, la lagrangiana trasforma come
\begin{equation}
    \Lag \to \Lag + a^\mu\p_\mu\Lag + \cdots = \Lag + a^\nu\p_\mu(\delta^\mu_\nu\Lag),
\end{equation}
dove si può identificare di nuovo $\epsilon \equiv a^\nu$ e $\p_\mu k^\mu \equiv \p_\mu(\delta^\mu_\nu\Lag)$. La corrente di Noether quindi è un tensore a due indici
\begin{equation}
    J^\mu_\nu = \frac{\p\Lag}{\p(\p_\mu\phi)}\p_\nu\phi - \delta^\mu_\nu\Lag,
\end{equation}
che, tramite la metrica, si può scrivere nella forma
\begin{equation}
    T^{\mu\nu} = J^\mu_\rho \eta^{\rho\nu},
\end{equation}
indicata con il nome di ``tensore energia-impulso'', soddisfacente l'equazione di conservazione \eqref{eq: legge di conservazione di Noether}. Le cariche conservate associate sono
\begin{equation}
    p^\mu = \int d^3x\,T^{0\mu},
\end{equation}
che altro non sono che le componenti del quadrimomento. Si potrebbe fare lo stesso per una trasformazione del gruppo di Lorentz, a cui corrisponde il tensore conservato
\begin{equation}
    M^\mu_{\rho\sigma} = -x_\rho T^\mu_\nu + x_\sigma T^\mu_\rho,
\end{equation}
e le sei cariche conservate
\begin{equation}
    L_{\rho\sigma} = \int d^3x\, M^0_{\rho\sigma}.
\end{equation}
L'insieme delle $p^\mu$ e delle $L_{\rho\sigma}$ contiene tutte le cariche conservate nello spaziotempo 4-dimensionale. Tuttavia, le trasformazioni introdotte più sopra non sono tutte le possibili simmetrie, ne esistono altre chiamate ``simmetrie interne''. Ad esempio, se si considera un campo scalare complesso avente lagrangiana
\begin{equation}
    \Lag = \p_\mu\phi\p^\mu\phi^* - m^2|\phi|^2,\quad |\phi|^2 \equiv \phi\phi^*,
\end{equation}
si vede che questa è invariante sotto trasformazioni di Poincaré, ma anche sotto trasformazioni del tipo $\phi \to \phi' \equiv e^{i\alpha}\phi$, con $\alpha \in \bb{R}$, appartenenti al gruppo $U(1)$. Espandendo ora i campi $\phi'$ in serie di Taylor, scrivendo così le trasformazioni
\begin{equation}
    \phi \to \phi + i\alpha\phi + \cdots,\quad \phi^* \to \phi^* - i\alpha\phi^* + \cdots,
\end{equation}
si può derivare la forma della corrente come
\begin{equation}
    J^\mu = i\phi\p^\mu\phi^* - i\phi^*\p^\mu\phi,
\end{equation}
dove si è definito $\p^\mu \equiv \eta^{\mu\nu}\p_\nu$ in maniera simile a quanto fatto per $x_\mu$. Il fatto che questa corrente soddisfi la legge di conservazione \eqref{eq: legge di conservazione di Noether} non è immediatamente ovvio, però, dal calcolo esplicito, si può trovare che $\p_\mu J^\mu = 0$ se e solo se il campo complesso $\phi$ soddisfa le equazioni del moto. 

\section{Geometria differenziale}
In generale, su una varietà $M$ si possono definire degli aperti $M_i \subseteq M$, dette carte, a cui possono essere associati degli ulteriori aperti $A_i \subseteq \bb{R}^D$, o in $\bb{R}^{1,D - 1}$ nel caso di una metrica minkowskiana, tramite delle mappe $\phi : M \to \bb{R}^D$ tali che $\phi_i(P) = (x^0,\dots,x^{D - 1})$, dove $P \in M_i \subseteq M$. Il punto può appartenere ad una sola carta oppure a più carte che in quel punto, ed eventualmente nel suo intorno, si sovrappongono. Nel caso in cui questo avvenga, è possibile definire, a partire dalle mappe $\phi_i$, delle mappe che legano due diversi aperti $A_i \subseteq \bb{R}^D$ tra di loro; tali mappe prendono il nome di cambi di coordinate. 

Tra gli oggetti che possono essere definiti su una varietà $M$ ci sono i campi vettoriali $V^\mu(x)$. Più precisamente, i campi vettoriali sono definiti come mappe che associano ad ogni punto $p$ della varietà un vettore appartenente allo spazio tangente in quel punto $T_pM$, ossia
\begin{equation}
    \forall p \in M,\quad V : M \to T_pM,\quad p \mapsto V(p) \in T_pM.
\end{equation}
In maniera simile si può definire la controparte ``covariante'' dei campi vettoriali, che in un linguaggio più moderno prendono il nome di 1-forme differenziali. Una 1-forma differenziale $\omega_\mu$ è una mappa che associa ad ogni punto $p$ della varietà un vettore appartenente allo spazio cotangente in quel punto $T_p^*M$ , ossia
\begin{equation}
    \forall p \in M,\quad \omega : M \to T_p^*M,\quad p \mapsto \omega(p) \in T_p^*M.
\end{equation}
In modo equivalente, una 1-forma si può definire come una mappa tra lo spazio tangente e un aperto di $\bb{R}$, per cui
\begin{equation}
    \forall v \in T_pM,\quad \omega : T_pM \to \bb{R},\quad v \mapsto \omega(v) \in \bb{R}.    
\end{equation}
Campi vettoriali e 1-forme differenziali trasformano sotto trasformazioni di coordinate $x^\mu \to x'^\mu$ con lo jacobiano e lo jacobiano inverso rispettivamente, per cui
\begin{equation}
    V'^\mu(x') = \frac{\p x'^\mu}{\p x^\nu}V^\nu(x),\quad \omega'_\mu(x') = \frac{\p x_\mu}{\p x'_\nu}\omega_\nu(x).
\end{equation}
In generale, esistono infinite basi con cui si posso esprimere campi vettoriali e 1-forme, per ora si è dato per si è sottinteso l'uso della base coordinata $\{\p_\mu\}$ con la quale si può scrivere 
\begin{equation}
    V = V^\mu\p_\mu = V'^\mu\p'_\mu.
\end{equation}
La base per le 1-forme è definita invece tramite la relazione
\begin{equation}
    \mybraket{dx^\mu}{\p_\nu} = \delta^\mu_\nu.
\end{equation}
Le 1-forme sono particolarmente importanti perché posso essere viste come campi dinamici sulla varietà. 

A questo punto si possono definire i tensori come oggetti del tipo $T^{(q,r)} \equiv T \in \ot^q T_p^*M \ot^r T_pM$, dove $(q,r)$ è il rango del tensore. In componenti, un tensore ha la forma
\begin{equation}
    T = T^{\mu_1\cdots\mu_r}_{\nu_1\cdots\nu_q}\frac{\p}{\p x^{\mu_1}}\cdots\frac{\p}{\p x^{\mu_r}}dx^{\nu_1}\cdots dx^{\nu_q}.
\end{equation}
Le forme differenziali sono tensori del tipo $T_{\nu_1\cdots\nu_q} dx^{\nu_1}\cdots dx^{\nu_q}$ che appartengono allo spazio cotangente, dove $T_{\nu_1\cdots\nu_q}$ ha gli indici anti-simmetrizzati, cioè tali che sotto uno scambio $\nu_i \leftrightarrow \nu_j$ si ha un cambio di segno. Più precisamente, si può definire 
\begin{equation}
    T_{\nu_1\cdots\nu_q} dx^{\nu_1}\wedge \cdots\wedge dx^{\nu_q},
\end{equation}
dove $dx^{\nu}\wedge dx^{\nu} = dx^\mu \ot dx^\nu - dx^\nu \ot dx^\mu$ con $\wedge$ il prodotto esterno. Tramite il prodotto $\wedge$ si può definire una $p$-forma differenziale come
\begin{equation}
    \omega = \frac{1}{p!}\omega_{\mu_1\cdots\mu_p}dx^{\mu_1}\wedge\cdots\wedge dx^{\mu_p} \in \Omega^p(M),
\end{equation}
dove $\Omega^p(M)$ è lo spazio delle $p$-forme differenziali definite sulla varietà $M$. Supponendo di avere due forme differenziali $\omega^{(p)}$ e $\omega^{(q)}$, di rango $p$ e $q$ rispettivamente, e definendo la derivata esterna $d$ come la mappa
\begin{equation}
    d : \Omega^p(M) \to \Omega^{p + 1}(M),\quad \omega^{(p)} \mapsto d\omega^{(p)} \in \Omega^{p + 1}(M),
\end{equation}
quindi la derivata esterna di una $p$-forma si può scrivere come
\begin{equation}
    d\omega^{(p)} = \frac{1}{p!}\p_\mu\omega_{\mu_1\cdots\mu_p}dx^\mu\wedge dx^{\mu_1}\wedge \cdots \wedge dx^{\mu_p}.
\end{equation}
si possono verificare le seguenti proprietà:
\begin{enumerate}
    \item il prodotto esterno di $\omega^{(p)}$ e $\omega^{(q)}$ restituisce una $(p + q)$-forma differenziale, ossia $\omega^{(p)}\wedge\omega^{(q)} = \omega^{(p + q)}$;
    \item la derivata esterna è nilpotente, ossia $dd(\cdots) = 0$;
    \item la derivata esterna del prodotto di due 1-forme è
    \begin{equation}
        d(\omega^{(p)}\wedge\omega^{(q)}) = d\omega^{(p)} \wedge \omega^{(q)} + (-1)^p\omega^{(p)} \wedge d\omega^{(q)} \in \Omega^{(p + q + 1)}(M).
    \end{equation}
\end{enumerate}
Si può definire poi la contrazione, o prodotto interno, $i_V$ come la mappa 
\begin{equation}
    i_V : \Omega^p(M) \to \Omega^{p - 1}(M),\quad \omega^{(p)} \mapsto i_V\omega^{(p)} \in \Omega^{p - 1}(M),\, V \in T_pM,
\end{equation}
quindi la contrazione di una $p$-forma si può scrivere come
\begin{equation}
    i_V\omega^{(p)} = V^\mu\omega_{\mu\mu_1\cdots\mu_{p - 1}}dx^{\mu_1}\wedge \cdots \wedge dx^{\mu_{p - 1}}.
\end{equation}
Data una varietà differenziabile, si definisce la derivata di Lie di un campo vettoriale come 
\begin{equation}
    \Lie_XY = \comm{X}{Y},\quad X,Y \in T_pM.
\end{equation}
Nel caso di una $p$-forma differenziale $\omega \in \Omega^p(M)$, la sua derivata di Lie è data dalla formula magica di Cartan
\begin{equation}
\label{eq: formula magica di Cartan}
    \Lie_V\omega = d(i_V\omega) + i_Vd\omega,\quad V \in T_pM.
\end{equation}

\subsection{WIP}
Ritornando adesso al tensore metrico $g = g_{\mu\nu} dx^\mu \ot dx^\nu$, tramite $g_{\mu\nu}$ è possibile passare da elementi dello spazio tangente ad elementi dello spazio cotangente tramite la relazione 
\begin{equation}
    g_{\mu\nu}T^\nu = T_\mu,\quad T^\nu \in T_pM,\,T_\mu \in T_p^*M.
\end{equation}
Di fatto quindi la metrica costituisce un isomorfismo tra $T_pM$ e $T_p^*M$. 

Si introduce ora il formalismo dei vielbein. In generale in questo formalismo, si sceglie un ricoprimento di aperti della varietà dello spaziotempo $M$ e, per ciascuno di questi aperti, una base locale costituita da un insieme di $n$ campi vettoriali linearmente indipendenti 
\begin{equation}
    e_a = e_a^\mu \partial_\mu,\quad a = 1,\ldots,n,
\end{equation}
che insieme generano lo spazio tangente $n$-dimensionale in ogni punto dell'insieme. Dualmente, un vielbein determina, ed è determinato da, un co-vielbein, costituito da un insieme di $n$ 1-forme linearmente indipendenti,
\begin{equation}
    e^a = e^a_\mu \, dx^\mu,
\end{equation}
tali che
\begin{equation}
    e^a(e_b) = e^a_\mu e_b^\mu = \delta^a_b.
\end{equation}
Un vielbein viene solitamente specificato tramite i suoi coefficienti $e^a_\mu$ rispetto a una base coordinata, nonostante la scelta di un insieme di coordinate locali $x^\mu$ non sia necessaria per la definizione di un vielbein. Per $n = 4$, ossia per uno spaziotempo 4-dimensionale, il tensore metrico può quindi essere espresso come
\begin{equation}
    g_{\mu\nu} = e^a_\mu\eta_{ab}e^b_\nu,\quad \eta_{ab} = \diag{(-1,1,1,1)}.
\end{equation}
Per quanto riguarda i campi vettoriali, questi si possono scrivere come
\begin{equation}
    V = V^\mu\p_\mu = V^ae_a,\quad V^a = e^a_\mu V^\mu.
\end{equation}
Le azioni si costruiscono tramite le forme di grado massimo, che per uno spazio $D$-dimensionale sono 
\begin{equation}
    \omega^{(D)} = \frac{1}{D!}\omega_{0\cdots D - 1}dx^0\wedge \cdots \wedge dx^{D - 1},
\end{equation}
il cui integrale sulla varietà è definito come
\begin{equation}
    \int_M\omega^{(D)} = \int_M\omega_{0\cdots D - 1}dx^0 \wedge \cdots \wedge dx^{D - 1}.
\end{equation}
La forma volume è 
\begin{equation}
    dV = \sqrt{|\det{g}|}dx^0\cdots dx^{D - 1} = e\,dx^0\cdots dx^{D - 1},\quad e = \det{e^a_\mu},
\end{equation}
la quale è usata per definire l'azione 
\begin{equation}
    S = \int dV\,\Lag = \int d^Dx\,e\Lag.
\end{equation}
Un tensore, o meglio una densità tensoriale, utile è quello di Levi-Civita, che in un
\begin{equation}
    \epsilon_{a_0\cdots a_{D-1}} = 
    \begin{cases}
        1,\quad a_0\cdots a_{D-1}\,\rem{e'\,una\,permutazione\,pari} \\
        -1,\quad a_0\cdots a_{D-1}\,\rem{e'\,una\,permutazione\,dispari} \\
        0,\quad \rem{altrimenti}
    \end{cases}.
\end{equation}
Si può definire poi
\begin{equation}
    \begin{aligned}
        \epsilon^{b_0\cdots b_{D - 1}} &= \eta^{a_0b_0}\cdots\eta^{a_{D - 1}b_{D - 1}}\epsilon_{a_0\cdots a_{D - 1}} \\
        &= \begin{cases}
            -1,\quad a_0\cdots a_{D-1}\,\rem{e'\,una\,permutazione\,pari} \\
            1,\quad a_0\cdots a_{D-1}\,\rem{e'\,una\,permutazione\,dispari} \\
            0,\quad \rem{altrimenti}
        \end{cases},
    \end{aligned}
\end{equation}
per cui in maniera simile si definisce
\begin{equation}
    \begin{aligned}
        \epsilon_{\mu_0\cdots \mu_{D - 1}} &= e^{-1}\epsilon_{a_0\cdots a_{D - 1}}e^{a_0\cdots a_{D - 1}}_{\mu_0\cdots \mu_{D - 1}} \\
        &= \begin{cases}
            -1,\quad a_0\cdots a_{D-1}\,\rem{e'\,una\,permutazione\,pari} \\
            1,\quad a_0\cdots a_{D-1}\,\rem{e'\,una\,permutazione\,dispari} \\
            0,\quad \rem{altrimenti}
        \end{cases}.
    \end{aligned}
\end{equation}
Si noti che anche se si sta lavorando in $D$ dimensioni si può comunque definire $\epsilon_{a_1a_2}^{\mu_3\mu_4}$. Si può dimostrare che
\begin{equation}
    \begin{aligned}
        e^0 \wedge \cdots \wedge e^{D - 1} &= \frac{1}{D!}\epsilon_{a_0\cdots a_{D - 1}}e^{a_0}\wedge\cdots\wedge e^{a_{D - 1}} \\
        &= \sqrt{|\det{g}|}d^Dx,
    \end{aligned}
\end{equation}
dove si può usare $e^a = e^a_{\mu_1}dx^{\mu_1}$. 

Su una varietà dotata di una metrica si può definire un'altra operazione. Il duale di Hodge $*$ è definito come la mappa
\begin{equation}
    * : \Omega^p(M) \to \Omega^{D - p}(M),
\end{equation}
tale che 
\begin{equation}
    *e^0 \wedge \cdots \wedge e^{D - 1} = \frac{1}{(D - p)!}\epsilon^{a_1\cdots a_p}_{b_1\cdots b_{D -p}}e^{b_1}\wedge \cdots \wedge e^{b_{D - p}}.
\end{equation}
Applicando il duale di Hodge due volte si ottiene la mappa 
\begin{equation}
    ** : \Omega^p(M) \to \Omega^p(M),
\end{equation}
tale che 
\begin{equation}
    **\omega = - (-1)^{p(D - p) + t}\omega,
\end{equation}
dove $t = 0$ se la metrica è euclidea, o $t = 1$ se è minkowskiana. 

Si ricorda poi il teorema di Stokes
\begin{equation}
    \int_{\Sigma_p}d\omega^{(p - 1)} = \int_{\p\Sigma}\omega^{(p - 1)},
\end{equation}
dove $\p\Sigma_p$ è il bordo della varietà. 

Come detto prima, le forme differenziali si possono vedere come campi dinamici, l'esempio principale di ciò è il campo elettromagnetico $F = dA$ con il quale si possono scrivere le equazione di Maxwell. L'azione dell'Elettromagnetismo è
\begin{equation}
    S = -\frac{1}{4}\int d^4x F_{\mu\nu}F^{\mu\nu} = -\frac{1}{2}\int *F\wedge F,
\end{equation}
dove la seconda espressione è valida non solo in quattro dimensioni, ma in dimensioni generiche. Integrando per parti questa azione e imponendo $\delta S = 0$, si ottengono due delle equazioni di Maxwell nella forma
\begin{equation}
    d(* F) = 0.
\end{equation}
Le altre due sono derivate dall'identità di bianchi e sono contenute in
\begin{equation}
    dF = 0.
\end{equation}
L'invarianza di gauge è mantenuta in quanto per $A \to A + \Lambda$ si ha
\begin{equation}
    F \to F' \equiv dA + d(d\Lambda) = dA = F.
\end{equation}
La teoria di Maxwell può poi essere estesa alle $p$-forme. Sia $F^{(p+1)}$ una $(p + 1)$-forma definita da
\begin{equation}
    F^{(p + 1)} = dA^{(p)},
\end{equation}
l'azione della teoria è 
\begin{equation}
    S = -\frac{1}{2}\int *F^{(p + 1)}\wedge F^{(p + 1)}.
\end{equation}
Imponendo $\delta S = 0$ e dalla relazione $ddA^{(p)} = 0$ si trovano rispettivamente
\begin{equation}
    d(*F^{(p + 1)}) = 0,\quad dF^{(p + 1)} = 0.
\end{equation}
Si può dimostrare che
\begin{equation}
    * \omega^{(p)}\wedge\omega^{(p)} = \frac{1}{p!}\omega^{\mu_1\cdots\mu_p}\omega_{\mu_1\cdots\mu_p}\sqrt{|\det{g}}\,d^Dx,
\end{equation}
da cui per $p = 0$ si trova la proprietà 
\begin{equation}
    * 1 = \sqrt{|\det{g}}\,d^Dx \equiv dV\quad \omega^{(0)} := 1,
\end{equation}
permettendo l'uguaglianza
\begin{equation}
    \int dV\,f = \int_M*f.
\end{equation}
Su una varietà si possono definire delle derivate direzionali
\begin{align}
    \nabla_\mu V^\nu &= \p_\mu V^\nu + \Gamma_{\mu\rho}^\nu V^\rho, \\
    \nabla_\mu\omega_\nu &= \p_\mu\omega_\nu - \Gamma^\rho_{\mu\nu}\omega_\rho.
\end{align}
Tale nozione si può generalizzare ai tensori, in particolare per il tensore metrico si impone che 
\begin{equation}
    \nabla g_{\mu\nu} = \p_\mu g_{\nu\rho} - \Gamma^\lambda_{\mu\nu}g_{\lambda\rho} - \Gamma^\lambda_{\mu\rho}g_{\nu\lambda} \equiv 0.
\end{equation}
Tuttavia, questa imposizione non fissa la connessione; infatti, si deve anche richiedere che $\Gamma^\mu_{\nu\lambda} = \Gamma^\mu_{\lambda\nu}$ per fissare $\nabla$ in modo univoco e avere
\begin{equation}
    \Gamma^\rho_{\mu\nu} \equiv \left\{ \begin{matrix} \mu \\ \alpha\ \beta \end{matrix} \right\} = \frac{1}{2}g^{\rho\lambda}(\p_\mu g_{\lambda\nu} + \p_\nu g_{\lambda\mu} - \p_\lambda g_{\mu\nu}).
\end{equation}
Nel caso non sia verificata quest'ultima condizione, in generale si ha
\begin{equation}
    \Gamma^\rho_{\mu\nu} = \left\{ \begin{matrix} \mu \\ \alpha\ \beta \end{matrix} \right\} - K^\rho_{\mu\nu},
\end{equation}
il quale è effettivamente un tensore a differenza dei simboli di Christoffel, dove $K^\rho_{\mu\nu}$ è il tensore di contorsione. Si può poi definire il tensore di torsione come 
\begin{equation}
    T^\rho_{\mu\nu} = \Gamma^\rho_{\mu\nu} - \Gamma^\rho_{\nu\mu}.
\end{equation}
Si definisce infine il tensore di Riemann $R^\rho_{\mu\nu\sigma}$, con il quale si può definire il tensore di Ricci $R_{\mu\nu} = R^\lambda_{\mu\nu\lambda}$ e lo scalare di Ricci $R = g^{\mu\nu}R_{\mu\nu}$.

Partendo da $e^a$, si può definire una connessione con $\nabla_\mu e^a = -\omega^a_{\mu b}e^b$, o in maniera equivalente
\begin{equation}
    \p_\mu e^a_\nu - \Gamma^\rho_{\mu\nu}e^a_\rho = -\omega^a_{\mu b}e^b_\nu,
\end{equation}
da cui 
\begin{equation}
    \omega^a_{\mu b} = e^a_\nu e^\lambda_b\Gamma^\nu_{\mu\lambda} - e^\lambda_b\p_\mu e^a_\lambda.
\end{equation}
Da ciò si vede subito che $\nabla_\mu e^a = 0$ (fieldbein postulate). Si può definire una 1-forma connessione
\begin{equation}
    \omega^a_b = \omega^a_{\mu b}dx^\mu,
\end{equation}
con la quale si scrive la prima formula di struttura di Cartan 
\begin{equation}
\label{eq: prima formula di struttura di Cartan}
    de^a + \omega^a_b \wedge e^b = T^a,
\end{equation}
dove $T^a$ è la 2-forma di torsione. Espandendo $T^a$, si ha
\begin{equation}
    \frac{1}{2}T^a_{\mu\nu}dx^\mu \wedge dx^\nu = T^a.
\end{equation}
Ora se
\begin{equation}
    de^a = \p_{[\mu,e^a_\nu]}dx^\mu \wedge dx^\nu,\quad \p_{[\mu,e^a_\nu]} \equiv \frac{1}{2}(\p_\mu e^a_\nu - \p_\nu e^a_\mu),
\end{equation}
allora 
\begin{equation}
    T^a_{\mu\nu} = \p_\mu e^a_\nu - \p_\nu e^a_\mu + \omega^a_{\mu\nu} - \omega^a_{\nu\mu},\quad \omega^a_{\mu\nu} = \omega^a_{\mu b}e^b_\nu.
\end{equation}
Si può quindi definire 
\begin{equation}
    T_{[\mu\nu]\rho} = T^a_{\mu\nu}e_{a\rho},\quad \omega_{\mu[\rho\nu]} = \omega^{ab}_\nu e_{a\rho}e_{b\nu},\quad \Omega_{[\mu\nu]\rho} = (\p_\mu e^a_\nu - \p_\nu e^a_\mu)e_{a\rho}.
\end{equation}
Così si ottiene la relazione
\begin{equation}
    T_{[\mu\nu]\rho} = \Omega_{[\mu\nu]\rho} + \omega_{\mu[\rho\nu]} - \omega_{\nu[\rho\mu]} + K_{\mu[\nu\rho]},
\end{equation}
da cui segue che 
\begin{equation}
    K_{\mu[\nu\rho]} = -\frac{1}{2}(T_{[\mu\nu]\rho} - T_{[\nu\rho]\mu} + T_{[\rho\mu]\nu}).
\end{equation}
Si può allora definire una derivata covariante lorentziana
\begin{equation}
    \begin{aligned}
        D_\mu V^a &= \p_\mu V^a + \omega^a_{\mu b} V^b, \\
        D_\mu V_a &= \p_\mu V_a - \omega_{\mu a}^b V_b,
    \end{aligned}
\end{equation}
in particolare per la metrica
\begin{equation}
    D_\mu\eta_{ab} = \p_\mu \eta_{ab} - \omega^c_{\mu a}\eta_{cb} - \omega^c_{\mu b}\eta_{ac} = -(\omega_{\mu ba} + \omega\mu ab) = 0,
\end{equation}
che conferma il risultato già noto dalla Relatività. 

Si ritorni adesso alla prima formula di Cartan \eqref{eq: prima formula di struttura di Cartan} e  si definisca la 2-forma 
\begin{equation}
    \rho^{ab} = \frac{1}{2}R^{ab}_{\mu\nu} dx^\mu \wedge dx^\nu.
\end{equation}
Ora dato che 
\begin{equation}
    \rho^{ab} = d\omega^a_b + \omega^a_c \wedge \omega^c_b,
\end{equation}
si ha
\begin{equation}
    R^{ab}_{\mu\nu} = \p_\mu\omega^{ab}_\nu - \p_\nu\omega^{ab}_\mu + \omega^{ac}_\mu\omega^b_{\nu c}- \omega^{ac}_\nu\omega^b_{\mu c},
\end{equation}
dove se al posto di $\omega$ si avesse avuto $\Gamma$ l'espressione sarebbe stata quella del tensore di Riemann, il quale comunque può essere definito come
\begin{equation}
    R^\rho_{\mu\nu\lambda} = R^{ab}_{\mu\nu}e_a^\rho e_{b\lambda},
\end{equation}
per cui 
\begin{equation}
    R = e^\mu_ae^\nu_bR_\mu^{ab}.
\end{equation}
Quindi, si può scrivere 
\begin{equation}
    \int dV\,R = \int d^Dx\,eR,
\end{equation}
pensando come campo indipendente il vielbein invece che la metrica. 

Infine, si definisce l'operatore aggiunto della derivata $d^\dag \equiv *d\,*$ come la mappa
\begin{equation}
    d^\dag : \Omega^p(M) \to \Omega^{p - 1}(M).
\end{equation}
Una proprietà interessante di questa mappa è
\begin{equation}
    *d*\omega^{(1)} = (-1)^D\nabla^\mu\omega_\mu,
\end{equation}
che per una $p$-forma diventa
\begin{equation}
    *d*\omega^{(p)} = (-1)^{t + 1 + (p - 1)(D - p)}\frac{1}{(p - 1)!}\nabla^\nu\omega_{\nu\mu_1\cdots\mu_{p - 1}}dx^{\mu_1}\wedge \cdots \wedge dx^{\mu_{p - 1}}.
\end{equation}
Si tenta ora di calcolare l'integrale 
\begin{equation}
    \int d^Dx\sqrt{-\det{g}}\nabla^\mu V_\mu.
\end{equation}
Per farlo si ricorda che 
\begin{equation}
    \nabla^\mu V_\mu = \delta^\mu_\nu\nabla_\mu V^\nu = \p_\mu V^\mu + \Gamma^\mu_{\mu\rho}V_\rho,
\end{equation}
dove
\begin{equation}
    \Gamma^\mu_{\mu\rho} = \frac{1}{2}g^{\mu\lambda}\p_\rho g_{\mu\lambda} = \frac{1}{\sqrt{-\det{g}}}\p_\rho\sqrt{-\det{g}}.
\end{equation}
Usando $\det{M} = \exp{(\tr{(M)})}$ con $M \equiv g_{\mu\nu}$, si ha 
\begin{equation}
    \frac{\p_\mu\det{M}}{\det{M}} = \tr{(M^{-1}\p_\mu M)},
\end{equation}
da cui
\begin{equation}
    \int d^Dx\sqrt{-\det{g}}\nabla^\mu V_\mu = \int d^Dx\bigg[\p_\mu(\sqrt{-\det{g}}V^\mu) - T^\mu_{\mu\nu}V^\nu\bigg],
\end{equation}
dove si è usato il fatto che per una qualsiasi connessione si ha $\Gamma^\mu_{\rho\nu} = \left\{ \begin{matrix} \mu \\ \alpha\ \beta \end{matrix} \right\} - K^\mu_{\rho\nu}$. 

\section{Spinori}
Ci si concentra ora sullo spazio di Minkowski 4-dimensionale e sul gruppo di Lorentz $\group{SO}{1,3}$. Se si guarda l'algebra associata $\alg{so}{1,3}$, si può dimostrare che questa è isomorfa alla somma diretta di due copie dell'algebra $\alg{so}{2}$ associata al gruppo $\group{SO}{2}$, ossia
\begin{equation}
    \alg{so}{1,3} \cong \alg{so}{2} \op \alg{so}{2}.
\end{equation}
I generatori dell'algebra si indicano tipicamente con $T_i$ e soddisfano la relazione di commutazione 
\begin{equation}
    \comm{T_i}{T_j} = if_{ijk}T_k,
\end{equation}
dove le $f_{ijk}$ sono le costanti di strutture, che in questo caso coincidono con il simbolo di Levi-Civita $\epsilon_{ijk}$; quindi, più precisamente
\begin{equation}
    \comm{T_i}{T_j} = i\epsilon_{ijk}T_k.
\end{equation}
Una possibile base dell'algebra di Lorentz si può scrivere tramite le matrici di Pauli $\sigma_i$, con le quali 
\begin{equation}
    T_i = \frac{\sigma_i}{2},
\end{equation}
dove 
\begin{equation}
    \sigma_1 = 
    \begin{pmatrix}
        0 & 1 \\
        1 & 0
    \end{pmatrix},\quad
    \sigma_2 = 
    \begin{pmatrix}
        0 & -i \\
        i & 0
    \end{pmatrix},\quad
    \sigma_3 = 
    \begin{pmatrix}
        1 & 0 \\
        0 & -1
    \end{pmatrix}.
\end{equation}
Le rappresentazioni di questo gruppo sono poi etichettate da un numero intero, o eventualmente semi-intero, indicato da $j \in \bb{Z}/2$ con cui si scrive la dimensione della rappresentazione $\rho$ data da
\begin{equation}
    \dim{\rho} = 2j + 1.
\end{equation}
Tuttavia, dato l'algebra di $\group{SO}{1,3}$ è somma diretta di due algebre, saranno necessari due numeri $j$ e $j'$ per identificare una rappresentazione del gruppo. In particolare, le rappresentazioni di dimensione minore inequivalenti sono due: $\bm{(1/2,0)}$ e $\bm{(0,1/2)}$; entrambe rappresentazioni 2-dimensionali. In generale, una rappresentazione si dice spinoriale se $j + j' \in \bb{Z}/2$. 

Il prossimo passo è notare che esiste un omomorfismo tra $\group{SL}{2,\bb{C}}$ e $\group{SO}{1,3}$. Si definiscono quindi le matrici 
\begin{equation}
    \sigma_\mu = (-\IM,\sigma_i),\quad \b{\sigma}_\mu = (\IM,\sigma_i),
\end{equation}
che soddisfano le proprietà:
\begin{enumerate}
    \item $\sigma_\mu\b{\sigma}_\nu + \b{\sigma}_\mu\sigma_\nu = 2\eta_{\mu\nu}\IM$;
    \item $\tr{(\sigma^\mu\b{\sigma}_\nu)} = 2\delta^\mu_\nu$.
\end{enumerate}
Ad ogni vettore, ad esempio alle coordinate $x^\mu$, si può associare una matrice $X$ definita come
\begin{equation}
    x^\mu \mapsto X := \sigma_\mu x^\mu = 
    \begin{pmatrix}
        x^0 + x^3 & x^1 - ix^2 \\
        x^1 + ix^2 & x^0 - x^3
    \end{pmatrix}.
\end{equation}
Si vede poi che il determinante è $\det{X} = -x_\mu x^\mu$ e si può dimostrare che la matrice $X$ è hermitiana, ovvero che  $X = X^\dag$. Usando proprio quest'ultima proprietà, si può dimostrare anche che 
\begin{equation}
    x^\mu = \frac{1}{2}\tr{(\sigma^\mu X)}.
\end{equation}
Si consideri quindi una matrice $A \in \group{SL}{2,\bb{C}}$, con questa si può definire la trasformazione 
\begin{equation}
    X \to X' = AXA^\dag : X'^\dag = X^\dag,\,\det{X} = \det{X'}.
\end{equation}
Questa è a tutti gli effetti una trasformazione di Lorentz. Si può quindi arrivare all'espressione dell'omomorfismo tra i gruppi $\group{SL}{2,\bb{C}}$ e $\group{SO}{1,3}$ di cui si parlava inizialmente
\begin{equation}
    (\Lambda^{-1})^\mu_\nu(A) = \frac{1}{2}\tr{(\sigma^\mu A\bb{\sigma}_\nu A^\dag)}.
\end{equation}
Il fatto che sia un omomorfismo si può codificare con la relazione $\group{SO}{1,3} \cong \group{SL}{2,\bb{C}}/\bb{Z}_2$, con ciò si parla di $\group{SL}{2,\bb{C}}$ come il \emph{doppio ricoprimento} di $\group{SO}{1,3}$. Tipicamente, si indica il gruppo $\group{SO}{1,3}$ tramite la nozione di gruppo di spin, indicandolo con $\group{Spin}{1,3}$. In generale, in $D$ dimensioni $\group{Spin}{1,D}$ è il ricoprimento di $\group{SO}{1,D}$ se la segnatura è quella di Minkowski, mentre $\group{Spin}{D}$ è il ricoprimento di $\group{SO}{D}$ se la segnatura è quella di Euclide. 

Si torni adesso alle rappresentazioni $\bm{(1/2,0)}$ e $\bm{(0,1/2)}$. Usando le relazioni
\begin{equation}
    A\b{\sigma}_\mu A^\dag = (\Lambda^{-1})^\mu_\nu\b{\sigma}_\nu,\quad A^\dag \sigma_\mu A = (\Lambda^{-1})^\nu_\mu\sigma_\nu,
\end{equation}
si può dimostrare che i generatori dell'algebra si possono scrivere come 
\begin{equation}
    \begin{cases}
        \sigma_{\mu\nu} = \dfrac{1}{4}(\sigma_\mu\b{\sigma}_\nu - \sigma_\nu\b{\sigma}_\mu) \\[2ex]
        \b{\sigma}_{\mu\nu} = \dfrac{1}{4}(\b{\sigma}_\mu\sigma_\nu - \b{\sigma}_\nu\sigma_\mu)
    \end{cases},
\end{equation}
con i quali si definiscono
\begin{equation}
    L(\lambda) := e^{\lambda^{\mu\nu}\sigma_{\mu\nu}/2},\quad \tilde{L}(\lambda) := e^{\lambda^{\mu\nu}\b{\sigma}_{\mu\nu}/2},
\end{equation}
dove $\lambda^{\mu\nu} = -\lambda^{\nu\mu}$ e, dato che $\sigma_{\mu\nu}^\dag = -\b{\sigma}_{\mu\nu}$, sono tali che
\begin{equation}
    L(\lambda)^\dag = \tilde{L}(\lambda)^{-1}.
\end{equation}
Le $L(\lambda)$ e le $\tilde{L}(\lambda)$ consistono in una parametrizzazione delle matrici $A \in \group{SL}{2,\bb{C}}$. Si è così costruita la struttura che si desiderava: due parti che non trasformano simultaneamente associate a $\bm{(1/2,0)}$ per $L(\lambda)$ e a $\bm{(0,1/2)}$ per $\tilde{L}(\lambda)$. Sotto una trasformazione di Lorentz $x^\mu \to \Lambda^\mu_\nu x^\nu$ si hanno
\begin{equation}
    \Psi(x) \to \Psi'(x') \equiv L^{-1}(\lambda)\Psi(\Lambda x),
\end{equation}
\begin{equation}
    \chi(x) \to \chi'(x') \equiv \tilde{L}^{-1}(\lambda)\chi(\Lambda x).
\end{equation}
Gli spinori $\Psi$ e $\chi$ sono tipicamente indicati come \emph{spinori left} e \emph{spinori right} rispettivamente, mentre globalmente prendono il nome di \emph{spinori di Weyl}, o \emph{spinori chirali}.

Si vogliono ora introdurre gli spinori di Dirac. Questi spinori esistono per qualsiasi dimensione dello spaziotempo, e permettono l'imposizione di condizioni di realtà. A tal proposito, in $D = 4$, si introducono matrici di Dirac
\begin{equation}
    \gamma^\mu := 
    \begin{pmatrix}
        \OM_2 & \sigma^\mu \\
        \b{\sigma}^\mu & \OM_2
    \end{pmatrix},\quad \dim{\gamma^\mu} = 4,
\end{equation}
tali che 
\begin{equation}
    \gamma^\mu\gamma^\nu + \gamma^\nu\gamma^\mu = 2\eta^{\mu\nu}\IM_4.
\end{equation}
Si possono poi definire le matrici 
\begin{equation}
    \gamma'^\mu := S^{-1}\gamma^\mu S,\quad S \in \group{GL}{4,\bb{C}},
\end{equation}
per le quali continua a valere la relazione precedente. Questa particolare relazione si anti-commutatore e si indica con
\begin{equation}
    \acomm{\gamma^\mu}{\gamma^\nu} = 2\eta^{\mu\nu}\IM_4,
\end{equation}
la quale, oltretutto, definisce l'algebra di Clifford per uno spaziotempo 4-dimensionale, indicata con $\alg{Cliff}{1,3}$. La definizione è poi estendibile a qualsiasi dimensione $D$ dello spaziotempo. Si definisce quindi la matrice
\begin{equation}
    \Sigma^{\mu\nu} := \frac{1}{4}\comm{\gamma^\mu}{\gamma^\nu},
\end{equation}
la quale soddisfa l'algebra di Lorentz, per cui
\begin{equation}
    L(\lambda) := e^{\lambda^{\mu\nu}\Sigma_{\mu\nu}} \in \group{Spin}{1,D - 1}.
\end{equation}
Le matrici $\gamma^\mu$ trasformano come un vettore di Lorentz 
\begin{equation}
    \gamma^\mu \to L(\lambda)^{-1}\gamma^\mu L(\lambda) = \Lambda(\lambda)^\mu_\nu \gamma^\nu.
\end{equation}
La dimensione delle rappresentazioni dell'algebra di Clifford è $2^{D/2} \times 2^{D/2}$. Dunque, si definisce lo spinore di Dirac 
\begin{equation}
    \Psi(x) \in \bb{C}^{2^{D/2}},\quad \Psi(x) \to \Psi'(x') = L(\lambda)^{-1}\Psi(\Lambda x).
\end{equation}
\begin{ex}(rotazione di uno spinore)
    Considerando la particolare trasformazione $\lambda^{12} \equiv \theta \neq 0$, si ha 
    \begin{equation}
        L(\lambda) = e^{\theta\Sigma_{12}} = 1 + \theta\Sigma_{12},
    \end{equation}
    dove le potenze pari sono proporzionali all'identità; infatti
    \begin{equation}
        \Sigma_{12}\Sigma_{12} = \frac{1}{2}(\gamma^1\gamma^2)\frac{1}{2}(\gamma^1\gamma^2) = -\frac{1}{4}\IM_4.
    \end{equation}
    Quindi, si ottiene la relazione
    \begin{equation}
        \cos{\bigg(\frac{\theta}{2}\bigg)}\IM + 2\sin{\bigg(\frac{\theta}{2}\bigg)}\Sigma_{12},
    \end{equation}
    da cui si vede che per $\theta = 2\pi$, si ottiene $-\IM$. 
\end{ex}
Si torni in $D = 4$. Qui si ha 
\begin{equation}
    \Sigma^{\mu\nu} = 
    \begin{pmatrix}
        \sigma_{\mu\nu} & \OM_2 \\
        \OM_2 & \b{\sigma}_{\mu\nu}
    \end{pmatrix},
\end{equation}
andando ad agire su $\Psi(x)$ le due componenti sono indipendenti,
\begin{equation}
    \Psi(x) = 
    \begin{pmatrix}
        \Psi(x) \\
        \chi(x)
    \end{pmatrix},
\end{equation}
dove la componente left trasforma con $\sigma_{\mu\nu}$, mentre quella right con $\b{\sigma}_{\mu\nu}$. Tutto cioò è ripetibile in $D$ dimensioni, ottnenendo la relazione
\begin{equation}
    \underbrace{2^{D/2}}_{rappr.\,riducibile} = \underbrace{2^{D/2 - 1}}_{irrep} + \underbrace{2^{D/2 - 1}}_{irrep}.
\end{equation}

Al campo spinoriale si deve poi dare una dinamica. L'equazion del moto è l'equazione di Dirac
\begin{equation}
\label{eq: equazione di Dirac}
    (\sl{\p} - m)\Psi = 0,\quad \sl{\p} \equiv \gamma^\mu\p_\mu.
\end{equation}
Per ricavare l'equazione \eqref{eq: equazione di Dirac}, si deve definire un'azione per il campo spinoriale. In realtà, la scelta è molto limitata dalle condizioni:
\begin{enumerate}
    \item l'azione deve essere lineare nelle derivate;
    \item l'azione deve essere reale;
    \item l'azione deve essere un invariante di Lorentz.
\end{enumerate}
Un termine invariante lineare nelle derivate è $\gamma^\mu\p_\mu$, ma si deve introdurre anche il campo $\Psi$. Per farlo, si nota che su $\bb{C}^4$ esiste un prodotto interno naturale dato da
\begin{equation}
    \Psi^\dag\Psi \in \bb{R}.
\end{equation}
Essendo uno scalare, si potrebbe pensare che sia anche invariante di Lorentz; tuttavia, usando le trasformazioni
\begin{equation}
    \begin{cases}
        \Psi \to \Psi' \equiv L^{-1}(\lambda)\Psi(\Lambda x) \\
        \Psi^\dag \to \Psi'^\dag \equiv \Psi^\dag(\Lambda x)(L^{-1}(\lambda))^\dag
    \end{cases},
\end{equation}
si trova che $\Psi^\dag\Psi$ non è un invariante di Lorentz. Usando però il fatto che, in generale, $\b{\Psi} \equiv \Psi^\dag\beta$ con $\beta \gamma^\mu \beta^{-1} = -(\gamma^\mu)^\dag$ è un invariante di Lorentz, si può costruire l'invariante 
\begin{equation}
    \b{\Psi} = \Psi^\dag i\gamma^0\Psi,
\end{equation}
dove si è scelto $\beta \equiv i\gamma^0$. L'azione ha dunque la forma
\begin{equation}
    S[\Psi,\b{\Psi}] = \int d^4x\,(m\b{\Psi}\Psi - \b{\Psi}\sl{\p}\Psi).
\end{equation}
Dal principio di minima azione applicato alla variazione 
\begin{equation}
    \delta S = \int d^4x\,[\delta\b{\Psi}(\sl{\p} - m)\Psi - (\p_\mu\b{\Psi}\gamma^\mu + \b{\Psi}m)\delta\Psi],
\end{equation}
si ottengono le due equazioni del moto equivalenti
\begin{equation}
    \begin{cases}
        (\sl{\p} - m)\Psi = 0 \\
        \b{\Psi}(\overset{\leftarrow}{\sl{\p}} - m) = 0
    \end{cases},\quad \overset{\leftarrow}{\sl{\p}} \equiv \p_\mu\b{\Psi}\gamma^\mu.
\end{equation}
Si può infine dimostrare perché la dimensione della rappresentazione sia $2^{D/2}$. Si parte da $D = 2$, si devono trovare due matrici $\gamma^1$ e $\gamma^2$ tali che $\acomm{\gamma^1}{\gamma^2} = 0$ e che $(\gamma^1)^2 = (\gamma^2)^2 = \IM_4$. Una scelta possibile è
\begin{equation}
    \gamma^1 = \sigma_1,\quad \gamma^2 = \sigma_2,
\end{equation}
definendo $\alg{Cliff}{2}$. Per $D = 3$ si aggiunge $\gamma^3 = \sigma_3$, definendo $\alg{Cliff}{3}$. Per $D = 4$ però si deve ricorrere al prodotto tensoriale, o di Kronecker, definito da
\begin{equation}
    A \ot B = 
    \begin{pmatrix}
        A_{11}B & \cdots & A_{1n}B \\
        \vdots & \ddots & \vdots \\
        A_{n1}B & \cdots & A_{nn}B
    \end{pmatrix},\quad \dim{(A \ot B)} = nm \times nm.
\end{equation}
Quindi, si scelgono
\begin{equation}
    \begin{cases}
        \gamma^0 = \sigma_1 \ot \IM \\
        \gamma^1 = \sigma_2 \ot \IM \\
        \gamma^2 = \sigma_3 \ot \sigma_1 \\
        \gamma^3 = \sigma_3 \ot \sigma_2
    \end{cases}.
\end{equation}
Se $D = 2m$, allora $m$ corrisponde al numero di matrici che compaiono in ogni prodotto tensoriale, dove per costruire le matrici successive si usa un metodo iterativo che consiste nel costruire il prodotto tensoriale tra $\sigma_3$ e le matrici precedenti, ossia
\begin{equation}
    \begin{cases}
        \gamma^1 = \sigma_1 \ot \IM \ot \cdots \ot \IM \\
        \gamma^2 = \sigma_2 \ot \IM \ot \cdots \ot \IM \\
        \gamma^3 = \sigma_3 \ot \sigma_1 \ot \IM \ot \cdots \ot \IM \\
        \gamma^4 = \sigma_3 \ot \sigma_2 \ot \IM \ot \cdots \ot \IM \\
        \gamma^5 = \sigma^3 \ot \gamma^3 \\
        \gamma^6 = \sigma^3 \ot \gamma^4 \\
        \vdots
    \end{cases}.
\end{equation}
Dunque, la dimensione è $2^m = 2^{D/2}$. L'algebra di Clifford esiste sicuramente ed si costruisce in questo modo in dimensione pari. 

Per costruire le matrici $\gamma^\mu$ in dimensione dispari $D = 2m + 1$; ad esempio, in $D = 7$ si definisce 
\begin{equation}
    \gamma^7  = \sigma_3 \ot \sigma_3 \ot \sigma_3.
\end{equation}
La dimensione continua ad essere $2^m \times 2^m$. 

Si può passare dalla rappresentazione di Euclide a quella Lorentziana con $\gamma^0 \equiv i\gamma^0$ tale che 
\begin{equation}
    (\gamma^0)^2 = -\IM,\quad (\gamma^0)^\dag = -\gamma^0,
\end{equation}
oppure più in generale che 
\begin{equation}
    (\gamma^\mu)^\dag = \gamma^0\gamma^\mu\gamma^0. 
\end{equation}
In dimensioni pari esiste una sola rappresentazione irriducibile, mentre in dimensione dispari ne esistono due ma portano alla stessa algebra. L'algebra contienr $\{\IM,\gamma^\mu,\gamma^\mu\gamma^\nu,\gamma^\mu\gamma^\nu\gamma^\rho\}$, ma in realtà la definzione di algebra di Clifford oermette di restringersi ad un insieme più piccolo; infatti
\begin{equation}
    \begin{split}
        \gamma^\mu\gamma^\nu &= \frac{1}{2}(\gamma^\mu\gamma^\nu + \gamma^\nu\gamma^\mu) + \frac{1}{2}(\gamma^\mu\gamma^\nu - \gamma^\nu\gamma^\mu) \\
        &= \eta^{\mu\nu} + \frac{1}{2}(\gamma^\mu\gamma^\nu - \gamma^\nu\gamma^\mu) \\
        &\equiv \eta^{\mu\nu} + \gamma^{[\mu}\gamma^{\nu]} \\
        &\equiv \eta^{\mu\nu} + \gamma^{\mu\nu},
    \end{split}
\end{equation}
dove si è definito
\begin{equation}
    \gamma^{[\mu}\gamma^{\nu]} \equiv \gamma^{\mu\nu} \equiv \frac{1}{2}(\gamma^\mu\gamma^\nu - \gamma^\nu\gamma^\mu).
\end{equation}
In generale quindi 
\begin{equation}
    \gamma^{\mu_1\cdots\mu_r} = \gamma^{[\mu_1}\gamma^{\mu_2}\cdots\gamma^{\mu_r]}.
\end{equation}
Si hanno inoltre 
\begin{equation}
    \gamma^\mu\gamma_\mu = D,
\end{equation}
\begin{equation}
    \gamma_\nu\gamma^\mu\gamma^\nu = \gamma_\nu(2\eta^{\mu\nu} - \gamma^\nu\gamma^mu) = (2 - D)\gamma^\mu,
\end{equation}
\begin{equation}
    \gamma^{\mu\nu}\gamma_\nu = ... = (D - 1)\gamma^\mu.
\end{equation}
Si definisce la matrice di chiralità
\begin{equation}
    \gamma^* := (-i)^{m + 1}\gamma^0\cdots \gamma^{D - 1},
\end{equation}
tale che 
\begin{equation}
    (\gamma^*)^2 = \IM,\quad \acomm{\gamma^*}{\gamma^\mu} = 0,\quad \tr{(\gamma^*)} = n_+ - n_- = 0,
\end{equation}
dove $n_+$ è il numero di autovalori positivi $+1$, mentre $n_-$ è quello di autovalori negativi $-1$, di conseguenza $n_+ = n_- \equiv n$ e
\begin{equation}
    n = 2^{m - 1}.
\end{equation}
Tramite la matrice di chiralità si può scrivere
\begin{equation}
    \gamma_{\mu_0\cdots\mu_D} = (i)^{m + 1}\epsilon_{\mu_0\cdots\mu_D}\gamma^*.
\end{equation}
La matrice di chiralità si può rappresentare come 
\begin{equation}
    \gamma^* = 
    \begin{pmatrix}
        \IM_{2^{m - 1}} & \OM \\
        \OM & \IM_{2^{m - 1}}
    \end{pmatrix}.
\end{equation}
Si definisce allora
\begin{equation}
    \gamma^\mu = 
    \begin{pmatrix}
        \alpha^\mu & \sigma^\mu \\
        \b{\sigma}^\mu & \beta^\mu
    \end{pmatrix},
\end{equation}
dove, per la relazione dell'algebra di Clifford, si ha $\alpha^\mu = \beta^\mu = 0$ e 
\begin{equation}
    \sigma^\mu\b{\sigma}^\nu + \sigma^\nu\b{\sigma}^\mu = \acomm{\sigma^\mu}{\b{\sigma}^\nu} = 2\eta^{\mu\nu},
\end{equation}
dove si definiscono le matrici 
\begin{equation}
    \sigma^\mu = (\IM,\sigma^i),\quad \b{\sigma}^\mu = (-\IM,\sigma^i),
\end{equation}
con $\sigma^i$ la generalizzazione delle matrici di Pauli, in quattro dimensioni $\acomm{\sigma^i}{\sigma^j} = \delta^{ij}$. Tutto ciò vale nella base di Weyl. Si definiscono poi i proiettori 
\begin{equation}
    P_L = \frac{1}{2}(\IM + \gamma^*),\quad P_R = \frac{1}{2}(\IM - \gamma^*),
\end{equation}
tali che 
\begin{equation}
    P_L^2 = P_L,\quad P_R^2 = P_R,\quad P_RP_L = P_LP_R = 0.
\end{equation}
Applicando il proiettore ad uno spinore di seleziona la componente corrispondente
\begin{equation}
    P_L\Psi = 
    \begin{pmatrix}
        \Psi \\
        0
    \end{pmatrix},\quad P_R\Psi = 
    \begin{pmatrix}
        0 \\
        \chi
    \end{pmatrix},
\end{equation}
relazioni che chiaramente non sono vere in ogni base ma solo in quella di Weyl. Si può dimostrare che
\begin{equation}
    \gamma^{\mu_1\cdots \mu_D}\gamma^* = -\frac{(-i)^{m + 1}}{(D - r)!}\epsilon^{\mu_r\cdots\mu_1\nu_1\cdots\nu_{D - r}}\gamma_{\nu_1\cdots \nu_{D - r}}.
\end{equation}
In $D = 2m + 1$,
\begin{equation}
    \gamma^\mu \in \{\gamma^0,\gamma^1,\dots,\gamma^{2m-1},\gamma^{2m} \equiv \gamma^*\}
\end{equation}
oppure 
\begin{equation}
    \gamma^\mu \in \{\gamma^0,\gamma^1,\dots,\gamma^{2m-1},\gamma^{2m} \equiv -\gamma^*\},
\end{equation}
che sono scelte non equivalenti, ma portano ad algebre equivalenti. 

\paragraph{Intertwining operators.} Se $\b{\Psi} = \Psi^\dag\gamma^0$, si definisce $A = i\gamma^0$ tale che 
\begin{equation}
    A\gamma^\mu A^{-1} = -(\gamma^\mu)^\dag \quad\Longleftrightarrow\quad A\gamma^\mu = -(\gamma^\mu)^\dag A,
\end{equation}
con il quale quindi si passa da $\gamma$ ad $\gamma^\dag$. Si definisce poi $B$ come la mappa che permette di passare da $\gamma$ a $\gamma^*$ e con $C$ quella tra $\gamma$ e $\gamma^T$, come 
\begin{equation}
    (\gamma^\mu)^* = -t_0t_1B\gamma^\mu B^{-1},
\end{equation}
\begin{equation}
    (\gamma^\mu)^T = +t_0t_1C\gamma^\mu C^{-1},
\end{equation}
dove $t_0 = t_1 = \pm 1$ e con $BB^\dag = CC^\dag = \IM$ e con $C^T = -t_0C$ o $C^* = -t_0C^{-1}$. Si definiscono quindi le $\Gamma^{(\nu)} \in \{\IM,\gamma^\mu,\gamma^{\mu\nu},\dots\}$, tale che 
\begin{equation}
    (C\Gamma^{(\nu)})^T = t_r{(C\Gamma^{(r)})},\quad t_r = \pm 1.
\end{equation}
Quindi, $t_r$ con $r > 2$, è fissato da $t_0$ e $t_1$. Ad esempio
\begin{equation}
    \begin{split}
        (C\gamma^{\mu\nu})^T &= \frac{1}{2}(C\gamma^\mu\gamma^\nu - C\gamma^\nu\gamma^\mu) \\
        &= \frac{1}{2}[(\gamma^\nu)^T(\gamma^\mu)^T - (\gamma^\mu)^T(\gamma^\nu)^T]C^T \\
        &= -\frac{t_0}{2}C(\gamma^\nu\gamma^\mu - \gamma^\mu\gamma^\nu) \\
        &= t_0C\gamma^{\mu\nu}.
    \end{split}
\end{equation}
Da qui si trova che $t_2 = -t_0$, si può poi trovare che $t_3 = -t_1$, $t_{r + 4} = t_r$ e così via. Si può poi costruire esplicitamente una rappresentazione di $C$ come
\begin{equation}
    C_+ = \sigma_1 \ot \sigma_2 \ot \sigma_1 \ot \sigma_2 \ot \cdots,\quad t_0t_1 = 1,
\end{equation}
o
\begin{equation}
    C_- = \sigma_2 \ot \sigma_1 \ot \sigma_2 \ot \sigma_1 \ot \cdots,\quad t_0t_1 = -1,
\end{equation}
a differenza di $A = i\gamma^0$ che non dipende dalla segnatura. In modo simile, si ha che 
\begin{equation}
    B = it_0C\gamma^0 = t_0CA.
\end{equation}
Si dimostra 
\begin{equation}
    (\gamma^\mu)^\dag = (\gamma^T)^* = \cdots = -\gamma^0\gamma^\mu\gamma^0 = -A\gamma^\mu A^{-1}.
\end{equation}
Si dimostra
\begin{equation}
    \begin{split}
        B^*B &= -it_0C^*\gamma^*it_0C\gamma^0 \\
        &= -t_0C^{-1}(\gamma^0)^*C\gamma^0 \\
        &= -t_0C^{-1}B\gamma^0B^{-1}C\gamma^0(-t_0t_1) \\
        &= t_1(C^{-1}B)\gamma^0(C^{-1}B)^{-1}\gamma^0 \\
        &= -t_1.
    \end{split}
\end{equation}
In $D = 4$, da $B = \gamma^0\gamma^1\gamma^3$ e $C = i\gamma^3\gamma^1$ so ha $t_0 = 1$, allora $B = B^*$ e $BB^* = B^2 = \cdots = \IM$ e $t_1 = -1$, scegliendo quindi 
\begin{equation}
    t_0 = 1,t_1 = -1,t_2 = -1,t_3 = 1,\dots,
\end{equation}
invece che
\begin{equation}
    t_0 = 1,t_1 = 1,t_2 = -1,t_3 = -1,\dots.
\end{equation}
In dimensione $D$ generica, si era definito lo spinore di Dirac $\b{\Psi}$ con il quale si possono costruire gli invarianti di Lorentz. Si definisce in particolare lo \emph{spinore di Majorana}
\begin{equation}
    \b{\lambda} := \lambda^TC,
\end{equation}
con il quale si possono costruire nuovi invarianti. Quindi, applicando $\b{\lambda}$ a $\chi$ si ha
\begin{equation}
    \b{\lambda}\chi = \lambda^TC\chi,
\end{equation}
che trasforma come
\begin{equation}
    \delta(\b{\lambda}\chi) = (\delta\b{\lambda})\chi + \b{\lambda}\delta(\chi),
\end{equation}
dove
\begin{equation}
    \delta\lambda^T = -\frac{1}{4}\chi^T(\gamma_{\mu\nu})^T\lambda^{\mu\nu},\quad \delta\chi = -\frac{1}{4}\lambda^{\mu\nu}\gamma_{\mu\nu}\chi,
\end{equation}
si ha 
\begin{equation}
    \delta(\b{\lambda}\chi) = -\frac{1}{4}\lambda^{\mu\nu}[\lambda^T(\gamma_{\mu\nu})^TC\chi + \lambda^TC\gamma_{\mu\nu}\chi] = 0,
\end{equation}
da cui si definiscono gli invarianti $\lambda^T(\gamma_{\mu\nu})^TC\chi$ e $\lambda^TC\gamma_{\mu\nu}\chi$.

\section{Variabili di Grassman}
Le variabili di Grassman sono definite tramite un prodotto di anticommutazione
\begin{equation}
    \theta\eta = -\eta\theta.
\end{equation}
Sono nilpotenti $\theta^2 = 0$, per consistenza serve introdurre
\begin{equation}
    (\theta\eta)^* = \eta^*\theta^* = -\theta^*\eta^*.
\end{equation}
Una funzione $f(\theta)$ è una funzione lineare del tipo $f(\theta) = a + b\theta$, dove per distinguere le variabili di Grassman dai numeri complessi ordinari, si chiamano questi \emph{c-numeri} dove la ``c'' sta per \emph{commutanti}. A meno che non sia specificato, gli spinori avranno sempre come componenti numeri di Grassman. 

Quindi, in $\b{\lambda}\gamma^\mu\chi$, si può scegliere al posto della matrice di Dirac una generica matrice $M \in \alg{Cliff}{1,D - 1}$. Si considera quindi il caso in cui 
\begin{equation}
    M = \gamma_{\mu_1\cdots\mu_r},\quad r \le D,
\end{equation}
pensandolo come un isomorfismo tra l'algebra di Clifford e i numeri di Grassman. Esiste un'identità
\begin{equation}
    \b{\lambda}\gamma_{\mu_1\cdots\nu_r}\chi = t_r\b{\chi}\gamma_{\mu_1\cdots\mu_r}\lambda.
\end{equation}
Questa si può verificare ricordando che $s^T = s$ con $S \in \bb{C}$, per cui 
\begin{equation}
    \begin{split}
        (\b{\lambda}\gamma_{\mu_1\cdots\mu_r}\chi)^T &= \b{\lambda}\gamma_{\mu_1\cdots\mu_r}\chi \\
        &= -\chi^T(\gamma_{\mu_1\cdots\mu_r})^TC^T\lambda,
    \end{split}
\end{equation}
dove il segno meno nasce dallo scambio effettivo tra $\chi$ e $\lambda$. Quindi, ricordando che $\Gamma(r) = \gamma_{\mu_1\cdots\mu_r}$, si ha 
\begin{equation}
    \b{\lambda}\gamma_{\mu_1\cdots\mu_r}\chi = t_0\chi^TC\gamma_{\mu_1\cdots\mu_r}\lambda.
\end{equation}
Questa identità è detta \emph{Majorana flip}. Se $\lambda = \chi$ e $t_r = -1$, allora la quantità $\b{\lambda}\gamma^{\mu\nu}\lambda$ si annulla in $D = 4$, ma in generale dipende dalla dimensione. 

Per cambiare di posto un numero di spinori maggiore di due sono necessarie delle accortezze. Per quattro spinori ad esempio, il caso più comune, si ha l'identità di Fierz: per $\lambda_1$, $\lambda_2$, $\lambda_3$ e $\lambda_4$, si ha 
\begin{equation}
    (\b{\lambda}_1\lambda_2)(\b{\lambda_3}\lambda_4) = \nsum_A\frac{1}{2^m}\b{\lambda_1}\Gamma^A\lambda_4\b{\lambda_3}\Gamma^A\lambda_2,
\end{equation}
dove
\begin{equation}
    \Gamma^A \in \{\IM,\gamma^\mu,\dots,\gamma^{\mu_1\cdots\mu_D}\} \in \alg{Cliff}{2m}
\end{equation}
oppure 
\begin{equation}
    \Gamma^A \in \{\IM,\gamma^\mu,\dots,\gamma^{\mu_1\cdots\mu_m}\} \in \alg{Cliff}{2m + 1}.
\end{equation}

\subsection{Coniugazione di carica}
Prima di introdurre gli spinori di Majorana è necessario prima definire l'operazione di coniugazione di carica. Sia quindi $\lambda$ uno spinore di Dirac, si definisce la coniugazione di carica come
\begin{equation}
    \lambda^c := B^{-1}\lambda^*,
\end{equation} 
la quale trasforma secondo Lorentz. Facendone il coniugato di Majorana si ottiene
\begin{equation}
    \b{\lambda^c} = (B^{-1}\lambda^*)^TC = \lambda^\dag B^{-T}C.
\end{equation}
Si può dimostrare che $B^{-T} = -t_1B^{-1}$, con la quale 
\begin{equation}
    \b{\lambda^c} = -t_0t_1\lambda^\dag(i\gamma^0).
\end{equation}
Quindi, si vede come il coniugato di Majorana sia legato a quello di Dirac, manifestazione del fatto che $B$ e $C$ sono legate. 

Sia ora $M \in \alg{Cliff}{D}$ e sia $M^c := B^{-1}M^*B$, per una matrice di Dirac $\gamma_\mu$ si ha in particolare che il coniugato di carica è
\begin{equation}
    (\gamma_\mu)^c = -t_0t_1\gamma_\mu.
\end{equation}
Per la gamma chirale in particolare 
\begin{equation}
    \gamma_*^c = (-1)^{(D/2) + 1}\gamma_*.
\end{equation}
Si può poi calcolare che 
\begin{equation}
    \begin{split}
        (\b{\chi}M\lambda)^c = (\chi^TCM\lambda)^* &= (\chi_a(CM)_{ab}\lambda_b)^* \\
        &= \lambda_b^*(CM)^*_{ab}\chi_a^* \\
        &= \chi_a^*(CM)^*_{ab}\lambda_b^* \\
        &= -\chi^TC^*M^*\lambda^* \\
        &= i\chi^\dag\gamma^0 M^c \lambda^c,
    \end{split}
\end{equation}
dove si sono usate la definizione di $M^c$, $C^* = -t_0C^{-1}$ e $B = it_0C\gamma^0$, per cui
\begin{equation}
    (\b{\chi}M\lambda)^c = -t_0t_1\b{\chi^c}M^c\lambda^c.
\end{equation}
Dunque, l'operazione di coniugazione di carica agisce come il complesso coniugato a meno di un segno. 

\section{Spinori di Majorana}
Gli spinori di Dirac in generale hanno componenti complesse. Si vuole capire se è possibile imporre delle confizioni affinché le entrate siano invece reali. Una prima condizione che si potrebbe pensare di imporre è
\begin{equation}
    f - f^* = 0,
\end{equation}
dove $f$ è una qualsiasi funzione, così da avere $\Psi = \Psi^*$. Tuttavia, così facendo lo spinore non soddisfa le trasformazioni di Lorentz. Uno spinore che invece soddisfa la condizione di realtà delle componenti e le trasformazioni di Lorentz è il coniugato di carica
\begin{equation}
    \Psi^c = B^{-1}\Psi^*,
\end{equation}
con il quale si può imporre che $\Psi = \Psi^c$, equivalente a chiedere $\Psi^* = B\Psi$. Infatti, da $\Psi^* = B\Psi$ si ha 
\begin{equation}
    \Psi = B^*\Psi^* = B^*B\Psi,
\end{equation}
da cui si deve richiedere che $B^*B = \IM$, ma allora $t_1 = -1$. Possibile solo in alcune dimensioni, $D = 2,3,4 \mod{8}$, dove si ha che il coniugato di Majorana corrisponde all'aggiunto di Dirac, per cui 
\begin{equation}
    \b{\Psi^c} = \Psi^\dag A.
\end{equation}
\begin{obs}
    Ci sono rappresentazioni dove le matrici $C$ sono completamente reali (\emph{really-real representations}). Quando non esistono gli spinori di Majorana, si può introdurre una condizione pseudo-majorana. 
\end{obs}
Solo in dimensione $D = 2\mod{8}$ esistono gli spinori di Majorana che sono anche spinori di Weyl, detti \emph{spinori di Majorana-Weyl}. Per $D = 4$, non si può avere oggetti di questo tipo, ma si può avere o uno o l'altro caso, sia quindi $\Psi$ uno spinore di Majorana, in tal caso
\begin{equation}
    \Psi = (P_L + P_R)\Psi = P_L\Psi + P_R\Psi,
\end{equation}
dove $P_L\Psi$ e $P_R\Psi$ sono singolarmente spinori di Weyl, anche se la loro somma è uno spinore di Majorana. Si può vedere in particolare che 
\begin{equation}
    (P_L\Psi)^c = P_L^c\Psi^c = P_L^c\Psi = \frac{1}{2}(\IM + \gamma_*^c)\Psi,
\end{equation}
dove 
\begin{equation}
    \gamma_*^c = B^{-1}\gamma_*^*B = -\gamma_*,
\end{equation}
per cui
\begin{equation}
    P_L^c = \frac{1}{2}(\IM - \gamma_*) \equiv P_R,
\end{equation}
dunque 
\begin{equation}
    (P_L\Psi)^c = P_R\Psi,
\end{equation}
e ovviamente
\begin{equation}
    P_L\Psi = (P_R\Psi)^c.
\end{equation}
Quindi, la coniugazione di carica permette di passare da uno spinore all'altro, confermando la natura di spinore di Weyl dei singoli spinori. 

L'azione di Dirac era stata definita come 
\begin{equation}
    S_{Dirac} = -\int d^4x\,\b{\Psi}(\sl{\p} - m)\Psi,
\end{equation}
ma se $\Psi$ è uno spinore di Majorana la forma dell'azione non cambia dato che coniugato di Majorana corrisponde all'aggiunto di Dirac; quindi semplicemente
\begin{equation}
    S_{Majorana} = -\frac{1}{2}\int d^4x\,\b{\Psi}(\sl{\p} - m)\Psi.
\end{equation}
Tuttavia, si deve controllare se gli invarianti in $S_{Majorana}$ sono effettivamente tali. Se si suppone che i campi siano commutati, il Majorana flip restituisce che 
\begin{equation}
    \b{\Psi}\Psi = -\b{\Psi}\Psi,
\end{equation}
rendendo quindi il termine di massa in $S_{Majorana}$ nullo. Anche se questo di per se non sarebbe un problema dato che basterebbe richiedere che $m = 0$, il fatto è che anche il termine cinetico si annulla se i campi sono commutanti; infatti
\begin{equation}
    \b{\Psi}\gamma^\mu\p_\mu\Psi = \p_\mu\b{\Psi}\gamma^\mu\Psi\quad\Longrightarrow\quad \b{\Psi}\gamma^\mu\p_\mu\Psi = \frac{1}{2}\p_\mu(\b{\Psi}\gamma\Psi),
\end{equation}
che essendo un termine di bordo non contribuisce all'azione. In sintesi, se i campi sono commutanti l'azione è nulla, invece se questi anticommutano, allora l'azione non si annulla e sono presenti sia il termine cinetico che quello di massa. Calcolando la variazione $\delta S$, si ha 
\begin{equation}
    \b{\Psi}\delta\Psi = \delta\b{\Psi}\Psi
\end{equation}
e una relazione simile con le $\gamma^\mu$ incluse, per cui 
\begin{equation}
    \delta S = -\int d^4x\,\delta\b{\Psi}(\gamma^\mu\p_\mu - m)\Psi,
\end{equation}
da cui si ricava che lo spinore di Majorana soddisfa l'equazione di Dirac
\begin{equation}
    (\sl{\p} - m)\Psi = 0.
\end{equation}
Inoltre, si può scrivere l'azione di Majorana in funzione degli spinori di Weyl, 
\begin{equation}
    \begin{split}
        S_{Majorana} &= -\frac{1}{2}\int d^4x\,\b{\Psi}(\sl{\p} - m)(P_L + P_R)\Psi \\
        &= -\frac{1}{2}\int d^4x\,\Psi^TC(P_L + P_R)(\sl{\p} - m)(P_L + P_R)\Psi \\
        &= -\frac{1}{2}\int d^4x\,[2\b{\Psi}\Psi\gamma^\mu\p_\mu P_L\Psi - m\b{\Psi}P_L\Psi - m\b{\Psi}P_R\Psi] \\
        &\equiv S[P_L\Psi,(P_L\Psi)^c].
    \end{split}
\end{equation}

\subsection{Simmetrie}
Sia $\Psi$ uno spinore di Dirac. Non è molto difficile dimostrare che se si applica la trasformazione 
\begin{equation}
    \Psi \to \Psi' \equiv e^{i\theta}\Psi,\quad \theta \in \bb{R},
\end{equation}
l'azione è invariante. Questa trasformazione è un esempio di simmetria interna globale, in particolare una simmetria del gruppo $U(1)$. Se invece lo spinore $\Psi$ è uno spinore di Majorana, la trasformazione precedente non vede l'azione rimanere invariante, si può però operare la trasformazione
\begin{equation}
    \Psi \to \Psi' \equiv e^{i\theta\gamma_*}\Psi,
\end{equation}
che conserva la condizione di Majorana,
\begin{equation}
    B^{-1}\Psi^* \to B^{-1}e^{-i\theta\gamma_*^*}\Psi^* = B^{-1}e^{-i\theta\gamma_*^*}BB^{-1}\Psi^*,
\end{equation}
se $BB^{-1} = \IM$.
% DA VEDERE
% da cui, se $BB^{-1} = \IM$, si ha
% \begin{equation}
%     B^{-1}e^{-i\theta\gamma_*^*}B = \IM.
% \end{equation}
Si vede che 
\begin{equation}
    \begin{split}
        B^{-1}e^{-i\theta\gamma_*^*} &= B^{-1}[\IM - i\theta\gamma_*^* + \frac{1}{2}(-i)^2\theta^2(\gamma_*^*)^2 + \cdots] \\
        &= e^{i\theta\gamma_*},
    \end{split}
\end{equation}
dove si è usato che $(\gamma_*^*)^2 = \gamma_*^*BB^{-1}\gamma_*^*$ e per ogni scambio di $\theta$ con $\gamma_*$ c'è un segno. Si può far vedere che lo spinore $\b{\Psi}$ invece trasforma come 
\begin{equation}
    \b{\Psi} \equiv \Psi^TC \to \Psi^Te^{i\theta\gamma_*^T}C = \Psi^TCC^{-1}e^{i\theta\gamma_*^T}C = \b{\Psi}e^{i\theta\gamma_*}.
\end{equation}
Per quanto riguarda $\b{\Psi}\Psi$ invece, si può far vedere che non è invariante sotto tale trasformazione dato che gli esponenziali non si eliminano. Ponendo $m = 0$, il fattore cinetico diventa invece
\begin{equation}
    \b{\Psi}\gamma^\mu\p_\mu\Psi \to \b{\Psi}e^{i\theta\gamma_*}\gamma^\mu e^{i\theta\gamma_*}\p_\mu\Psi,
\end{equation}
ma se $\acomm{\gamma_*}{\gamma^\mu} = 0$, allora espandendo gli esponenziali e scambiando la $\gamma_*$, si ha 
\begin{equation}
    \b{\Psi}\gamma^\mu\p_\mu\Psi \to \b{\Psi}e^{i\theta\gamma_*}e^{-i\theta\gamma_*}\sl{\p}\Psi = \b{\Psi}\sl{\p}\Psi,
\end{equation}
che risulta quindi essere invariante. Da qui si giustifica la necessità di porre $m = 0$ nell'azione di Majorana. La corrente conservata è
\begin{equation}
    J^\mu_* = \frac{1}{2}\b{\Psi}\gamma^\mu\gamma_*\Psi,
\end{equation}
la quale soddisfa 
\begin{equation}
    \p_\mu J^\mu_* = -im\b{\Psi}\gamma_*\Psi,
\end{equation}
che è nulla appunto se $m = 0$. 

\section{Campo di Rarita-Schwinger}
Il campo di Rarita-Schwinger è un campo avente un indice spinoriale e un indice di Lorentz, è chiamato per questo \emph{spinore-vettore} ed è denotato con
\begin{equation}
    \Psi_{\alpha\mu}(x) \equiv \Psi_\mu(x),
\end{equation}
dove si è soppresso l'indice spinoriale per semplicità. In dimensione $D$ arbitraria, il numero di componenti in linea di principio sono $2^{[D/2]}D$, dove con $[D/2]$ se ne intende la parte intera, che in generale sono complesse. L'azione di questo campo deve soddisfare le condizioni già viste per il campo di Dirac:
\begin{enumerate}
    \item la lagrangiana deve essere invariante;
    \item l'azione deve essere reale;
    \item deve essere di prim'ordine nelle derivate $\p_\mu$;
\end{enumerate}  
insieme ad un'altra condizione. Per capire quale, si parte dall'azione del campo vettoriale più semplice, quella di Maxwell
\begin{equation}
    S_{Maxwell} = -\frac{1}{4}\int d^Dx\,F_{\mu\nu}F^{\mu\nu},
\end{equation}
dove $A_\mu$ è invariante per trasformazioni $A_\mu \to A_\mu + \p_\mu\Lambda$ con $\Lambda$ una funzione scalare. In generale, i gradi di libertà che soddisfano le equazioni del moto sono detti \emph{on-shell}, altrimenti sono detti \emph{off-shell}. Si fissano i gradi di libertà del campo $A_\mu$ tramite il gauge-fixing, in questo caso tramite il gauge di Coulomb $\p^iA_i = 0$. Dall'equazione del moto $\p^\mu F_{\mu\nu} = 0$ si separano 
\begin{equation}
    \begin{cases}
        \p^\mu F_{\mu 0} = 0 \\
        \p^\mu F_{\mu i} = 0
    \end{cases},
\end{equation}
che esplicitamente diventano
\begin{equation}
    \begin{cases}
        \p^i(\p_i A_0 - \p_0 A_i) = 0 \\
        \p^0(\p_0 A_i - \p_i A_0) + \p^j(\p_j A_i - \p_i A_j) = 0
    \end{cases},
\end{equation}
dove, considerando il gauge di Coulomb, dalla prima si trova che $\nabla^2 A_0 = 0$ che fissa $A_0 = 0$, con questo risultato dalla seconda si trova 
\begin{equation}
    \p^0\p_0 A_i + \p^j\p_j A_i \equiv \square A_i = 0,
\end{equation}
che è l'equazione dell'onda piana per il campo $A_i$. Il numero di gradi di libertà effettivo è $D - 2$ dove le due condizioni sono date dal gauge-fixing e dal fatto che $\v{k} \perp \v{A}$, ossai che $\v{k} \cdot \v{A} = 0$. 

Per il campo di Rarita-Schwinger si definsce allora la trasformazione 
\begin{equation}
    \psi_\mu \to \psi_\mu + \p_\mu\epsilon(x),
\end{equation}
dove $\epsilon$ è un parametro locale spinoriale. La lagrangiana deve essere invariante sotto questa trasformazione e di conseguenza ciò si ripercuote sull'azione, determinando la condizione aggiuntiva prima accennata. Una possibile azione qunidi è
\begin{equation}
    S_{R-S} = -\int d^Dx\,\b{\Psi}_\mu\gamma^{\mu\nu\rho}\p_\nu\Psi_\rho.
\end{equation}
La sua invarianza si può verificare, se $\b{Psi}_\mu = \psi_\mu^TC$, si ha 
\begin{equation}
    \begin{split}
        \delta(\b{\Psi}\gamma^{\mu\nu\rho}\p_\nu\Psi_\rho) &= \p_\mu\b{\epsilon}\gamma^{\mu\nu\rho}\p_\nu\Psi_\rho + \p_\mu\b{\epsilon}\gamma^{\mu\nu\rho}\p_\nu\p_\rho\epsilon \\
        &= \p_\mu\b{\epsilon}\gamma^{\mu\nu\rho}\p_\nu\Psi_\rho \\
        &= \p_\mu(\b{\epsilon}\gamma^{\mu\nu\rho}\p_\nu\Psi_\rho),
    \end{split}
\end{equation}
dove si è usata l'antisimmetria di $\gamma^{\mu\nu\rho}$; essendo una derivata totale non contribuisce all'azione. La variazione dell'azione è
\begin{equation}
    \begin{split}
        \delta_{R-S} &= -\int d^Dx (\delta\b{\Psi}_\mu\gamma^{\mu\nu\rho}\p_\nu\Psi_\rho - \b{\Psi}\gamma^{\mu\nu\rho}\p_\nu\delta\Psi_\rho) \\
        &= -\int d^Dx (\delta\b{\Psi}_\mu\gamma^{\mu\nu\rho}\p_\nu\Psi_\rho -\p_\nu\b{\Psi}_\mu\gamma^{\mu\nu\rho}\delta\Psi_\rho),
    \end{split}
\end{equation}
dove si è integrato per parti in secondo termine per ottenre una derivata totale. Quindi, si ottengono le due equazioni del moto del tutto equivalenti
\begin{equation}
    \begin{cases}
        \gamma^{\mu\nu\rho}\p_\nu\Psi_\rho = 0 \\
        \p_\nu\b{\Psi}_\rho\gamma^{\mu\nu\rho} = 0
    \end{cases}.
\end{equation} 
